\chapter{Architecture at a glance}
\label{ch:arch}

\prog is a three-crate Rust workspace with strict, mechanically-enforced boundaries.

\begin{center}
\begin{tabular}{@{}>{\ttfamily}l p{9.2cm}@{}}
\toprule
\normalfont\textbf{Crate} & \textbf{Role} \\
\midrule
sparkplug-b & Pure, generic Sparkplug~B library (protobuf, \code{seq}/\code{bdSeq}/rebirth
lifecycle). Zero application dependency — publishable to crates.io. \\
smart-me-client & Pure smart-me REST client (auth, endpoints, deserialization). Isolates
the heavy \code{reqwest}/\code{serde} dependencies. \\
smartme-bridge & The application: canonical \code{Measurement}, the pure
\code{Fresh|Stale|Failed} state machine, the two-task async runtime, adapters, web UI,
config, and wiring. \\
\bottomrule
\end{tabular}
\end{center}

The core that decides quality/staleness never imports async or transport crates —
\textbf{“no truth is ever decided inside an \code{async fn}”} — an invariant enforced at
compile time by \code{tests/arch\_purity.rs}. Dependency direction is one-way
(\code{sparkplug-b}, \code{smart-me-client} $\rightarrow$ \code{smartme-bridge}), enforced
by the Cargo graph and \code{cargo-deny}.

\section{The two-task runtime}

The running bridge is two asynchronous tasks joined by a single channel. They are started
together and neither is meaningful alone, so there is no window in which one runs without
the other.

\begin{center}
\begin{tabular}{@{}>{\ttfamily}l p{9.0cm}@{}}
\toprule
\normalfont\textbf{Task} & \textbf{Responsibility} \\
\midrule
poll\_publish & Paces the poll loop, fetches from the smart-me API under a deadline,
\textbf{judges} the result through the pure state machine, and forwards the verdict. \\
mqtt\_driver & Owns the Sparkplug session: \code{bdSeq}, the boot order, births, deaths
and reconnects. It transports the verdict; it never forms one. \\
\bottomrule
\end{tabular}
\end{center}

The split is the invariant, not an implementation detail. The state machine lives entirely
in the poll task and is mechanically barred from the MQTT task, so no quality decision can
be taken by the code that talks to the broker.

\subsection{Data flow}

\begin{center}
\code{smart-me REST} $\rightarrow$ \code{poll\_publish} $\rightarrow$
\emph{pure judgement} $\rightarrow$ channel $\rightarrow$ \code{mqtt\_driver}
$\rightarrow$ broker $\rightarrow$ Ignition
\end{center}

A fetch that exceeds its deadline is abandoned and counted as a timeout rather than
awaited, because a wedged loop stops publishing \textsc{stale} and so begins lying by
omission. The loop records a heartbeat \emph{before} each network call, never after: a
heartbeat written afterwards would look healthy exactly when the bridge is stuck.

\subsection{Why the transport gets its own task}

The MQTT event loop is polled in a task of its own rather than as a branch of a
\code{select!}. Polling it in a \code{select!} is not cancellation-safe: a cancelled
connect can abandon a half-finished CONNECT \emph{after} the broker has registered the
will, which then fires against a node that never birthed.

\section{The data flow, and where it can fail}
\label{sec:dataflow}

\begin{figure}[htbp]
\centering
\begin{tikzpicture}[
  node distance=6mm,
  box/.style={draw, rounded corners=2pt, align=center, inner sep=4pt, minimum height=8mm,
              font=\small},
  ext/.style={box, fill=black!4},
  flow/.style={-{Latex[length=2mm]}, thick},
  fail/.style={font=\scriptsize\itshape, align=left, text=black!65},
]
  \node[ext] (meter)  {the meter\\\scriptsize(Kamstrup, Telstar)};
  \node[ext, right=10mm of meter] (cloud) {smart-me\\cloud};
  \node[box, right=10mm of cloud] (poll)  {poll task\\\scriptsize one per meter};
  \node[box, below=8mm of poll] (oracle) {oracles\\\scriptsize pure judgement};
  \node[box, right=10mm of poll] (pub)   {publisher\\\scriptsize Sparkplug B};
  \node[ext, below=8mm of pub] (broker) {broker};
  \node[ext, below=8mm of broker] (scada) {Ignition};

  \draw[flow] (meter) -- (cloud);
  \draw[flow] (cloud) -- (poll);
  \draw[flow] (poll)  -- (oracle);
  \draw[flow] (oracle.east) -- ++(6mm,0) |- (pub.west);
  \draw[flow] (pub)   -- (broker);
  \draw[flow] (broker) -- (scada);

  \node[fail, below=2mm of meter.south] {stops measuring\\$\rightarrow$ \textsc{stale},\\
    \emph{feed-not-advancing}};
  \node[fail, below=2mm of cloud.south] {401 $\rightarrow$ latched\\429 $\rightarrow$ waits\\
    5xx/timeout $\rightarrow$ \textsc{stale}};
  \node[fail, right=2mm of oracle.east, yshift=-7mm] {refuses a reading\\it cannot vouch for};
  \node[fail, right=3mm of broker.east] {unreachable $\rightarrow$\\traced drop, never\\
    a silent loss};
\end{tikzpicture}
\caption{From meter to SCADA, with what each link does when it fails. Nothing that fails is
allowed to look like something that worked.}
\label{fig:dataflow}
\end{figure}

\textbf{Read the figure by its failure captions rather than its arrows.} The arrows are what
every such diagram has; what distinguishes this bridge is the second row --- each link, when it
breaks, produces a \emph{marked} outcome rather than a missing one. A meter that stops
measuring produces a stale value carrying its cause, not a gap. A cloud that refuses the
credential latches instead of retrying quietly. A broker that goes away produces a counted drop
rather than a reading nobody knows was lost.

\begin{note}
\textbf{The oracles sit off the main line on purpose.} They are pure: no clock of their own, no
network, no \code{async}. Everything they need arrives as an argument, which is what makes a
freshness verdict reproducible from a test rather than dependent on when it ran. The rule is enforced
mechanically: an architecture test refuses any \code{tokio}, \code{axum}, \code{reqwest} or
\code{rumqttc} import inside \code{core/} and \code{domain/}, so no truth can be decided
inside an \code{async fn}.
\end{note}
